\documentclass[12pt, a4paper]{article}

% --- Pacotes Fundamentais ---
\usepackage[utf8]{inputenc}
\usepackage[T1]{fontenc}
\usepackage[brazil]{babel}
\usepackage{geometry}
\usepackage{titlesec}
\usepackage{fancyhdr}
\usepackage{graphicx}
\usepackage{xcolor}
\usepackage{booktabs}
\usepackage{longtable}
\usepackage{tabularx}
\usepackage{colortbl}
\usepackage{float}
\usepackage{parskip}
\usepackage{lmodern}
\usepackage{lastpage}
\usepackage{pgfplots}
\pgfplotsset{compat=1.18}

% --- Configurações de Layout ---
\geometry{left=2.5cm, right=2.5cm, top=4cm, bottom=2.5cm, headsep=1cm, headheight=60pt}

% --- Definição de Cores ---
\definecolor{primaryColor}{RGB}{0, 51, 102}   % Azul Marinho
\definecolor{secondaryText}{RGB}{80, 80, 80}  % Cinza Escuro
\definecolor{accentColor}{RGB}{178, 34, 34}   % Vermelho Tijolo (Para destacar Valores Altos)
\definecolor{lightGray}{RGB}{245, 245, 245}
\definecolor{tableHeader}{RGB}{230, 230, 250}

% --- Configuração de Cabeçalho e Rodapé ---
\pagestyle{fancy}
\fancyhf{}

% Cabeçalho
\fancyhead[C]{%
    \begin{tabularx}{\textwidth}{@{} X r @{}}
        \textcolor{primaryColor}{\textbf{\footnotesize STUDIO DENTAL - VERSÃO 2}} & \scriptsize \textit{Relatório de Contas a Pagar} \\
        \textcolor{secondaryText}{\scriptsize CNPJ: 46.125.234/0001-66} & {\scriptsize \textbf{Data Base: 08/01/2026}}
    \end{tabularx}%
}

% Rodapé
\fancyfoot[R]{\normalsize Página \textbf{\thepage} de \textbf{\pageref{LastPage}}}

% --- Linha separadora do CABEÇALHO ---
\renewcommand{\headrulewidth}{1.5pt}
\renewcommand{\headrule}{%
    \vspace{0.1cm}
    \hbox to\headwidth{\color{primaryColor}\leaders\hrule height \headrulewidth\hfill}%
}

% --- Linha separadora do RODAPÉ ---
\renewcommand{\footrulewidth}{1.5pt}
\renewcommand{\footrule}{%
    \color{primaryColor}%
    \hbox to\headwidth{\leaders\hrule height \footrulewidth\hfill}%
    \vspace{0.2cm}
}

% --- Formatação de Títulos ---
\titleformat{\section}
{\color{primaryColor}\normalfont\Large\bfseries}
{\thesection}{1em}{}[{\titlerule[0.8pt]}]

\titleformat{\subsection}
{\color{primaryColor}\normalfont\large\bfseries}
{\thesubsection}{1em}{}

% --- Macro para Moeda ---
\newcommand{\reais}[1]{R\$ #1}

% --- INÍCIO DO DOCUMENTO ---
\begin{document}

% --- CAPA ---
\begin{titlepage}
    \centering
    \vspace*{3cm}
    {\Huge \textbf{\textcolor{primaryColor}{RELATÓRIO TÉCNICO}}} \\
    \vspace{0.5cm}
    {\Huge \textbf{\textcolor{primaryColor}{DE CONTAS A PAGAR}}} \\
    \vspace{2cm}
    \rule{10cm}{1.5pt} \\
    \vspace{0.5cm}
    {\Large \textit{Análise de Fluxo de Caixa e Projeção de Obrigações}} \\
    \vspace{1cm}
    \colorbox{lightGray}{%
        \parbox{0.9\textwidth}{%
            \centering
            \vspace{0.3cm}
            {\Large \textbf{STUDIO DENTAL}} \\
            \vspace{0.2cm}
            {\large \textcolor{accentColor}{\textbf{VERSÃO 2 - PROJEÇÃO AJUSTADA}}} \\
            \vspace{0.3cm}
            {\large CNPJ: 46.125.234/0001-66} \\
            \vspace{0.3cm}
        }%
    }
    \vspace{4cm}

    \vfill
    {\small Relatório técnico-financeiro baseado em análise de dados contábeis com vencimento a partir de 08/01/2026.} \\
    \vspace{0.5cm}
    {\small Valores superiores a \textcolor{accentColor}{\textbf{R\$ 5.000,00}} destacados em vermelho para atenção prioritária.} \\
    \vspace{0.3cm}
    {\small \textit{Nota: Pagamentos com vencimento em 07/01/2026 já foram efetuados e não constam neste relatório.}} \\
    \vspace{0.3cm}
    {\small \textcolor{accentColor}{\textbf{VERSÃO 2:}} Projeção ajustada excluindo UNIMED GOIÂNIA, SIMONE HELENA e SIMPLES NACIONAL reduzido a 2/3.} \\
    \vspace{1cm}
    \textbf{Data Base:} 08 de Janeiro de 2026
\end{titlepage}

% --- SEÇÃO 1: RESUMO EXECUTIVO ---
\section{Resumo Executivo}

Este relatório apresenta a análise completa das obrigações financeiras (contas a pagar) da empresa \textbf{STUDIO DENTAL} com vencimento igual ou posterior a \textbf{08 de Janeiro de 2026}.

A análise foi baseada exclusivamente nos dados extraídos do sistema contábil e processados de forma técnica e objetiva.

\subsection{Indicadores Principais}

\begin{table}[H]
    \centering
    \large
    \renewcommand{\arraystretch}{1.5}
    \begin{tabularx}{0.9\textwidth}{X r}
        \toprule
        \rowcolor{tableHeader} \textbf{Métrica} & \textbf{Valor} \\
        \midrule
        \textbf{Valor Total Geral a Pagar} & \textbf{\reais{629.317,28}} \\
        \textbf{Quantidade Total de Títulos} & \textbf{186 títulos} \\
        \textbf{Títulos de Alta Prioridade (> R\$ 5.000)} & \textbf{68 títulos} \\
        \textbf{Quantidade de Credores} & \textbf{24 fornecedores} \\
        \textbf{Ticket Médio por Título} & \textbf{\reais{3.383,43}} \\
        \midrule
        \rowcolor{accentColor!10} \textbf{Maior Obrigação Individual} & \textcolor{accentColor}{\textbf{\reais{11.292,42}}} \\
        \bottomrule
    \end{tabularx}
    \caption*{\footnotesize Valores processados com base em lançamentos contábeis de crédito (passivos).}
\end{table}

\subsection{Observações Críticas}

\begin{itemize}
    \item \textbf{Concentração Bancária Acentuada:} O fornecedor \textbf{SANTANDER BRASIL} representa \textbf{80,9\%} do valor total a pagar (\reais{509.312,73}), evidenciando forte dependência de financiamento bancário.

    \item \textbf{Redução Significativa:} VERSÃO 2 apresenta redução de \textbf{R\$ 109.597,92 (14,8\%)} em relação à projeção original, resultado da exclusão de UNIMED GOIÂNIA e ajuste do SIMPLES NACIONAL.

    \item \textbf{Projeção 2026:} Ano de 2026 representa \textbf{36,2\%} das obrigações (\reais{227.817,22}), com média mensal de R\$ 18.985, substancialmente inferior à versão original.

    \item \textbf{Títulos Prioritários:} Dos 186 títulos, \textbf{68 (36,6\%)} são superiores a R\$ 5.000,00, concentrados principalmente em financiamentos bancários de longo prazo.
\end{itemize}

\newpage

% --- SEÇÃO 2: DETALHAMENTO POR CREDOR ---
\section{Detalhamento por Credor}

A tabela a seguir apresenta o ranking dos 26 credores ativos, ordenados por valor total devido. Valores em \textcolor{accentColor}{\textbf{vermelho}} indicam obrigações superiores a R\$ 5.000,00.

\begin{longtable}{>{\raggedright\arraybackslash}p{9cm} r}
    \toprule
    \rowcolor{tableHeader} \textbf{Fornecedor / Parceiro} & \textbf{Valor Total (R\$)} \\
    \midrule
    \endfirsthead

    \multicolumn{2}{c}{{\tablename\ \thetable{} -- continuação}} \\
    \toprule
    \rowcolor{tableHeader} \textbf{Fornecedor / Parceiro} & \textbf{Valor Total (R\$)} \\
    \midrule
    \endhead

    \midrule
    \multicolumn{2}{r}{\textit{Continua na próxima página...}} \\
    \endfoot

    \bottomrule
    \endlastfoot

    SANTANDER BRASIL & \textcolor{accentColor}{\textbf{509.312,73}} \\
    SIMPLES NACIONAL & \textcolor{accentColor}{\textbf{41.602,08}} \\
    VICTALAB F MANIPULACAO LTDA & \textcolor{accentColor}{\textbf{16.263,64}} \\
    MR TRIPLAN CONTABILIDADE E CONSULTORIA EIRELI & \textcolor{accentColor}{\textbf{14.400,00}} \\
    GFIP SEFIP 8,40 GRF FGTS & \textcolor{accentColor}{\textbf{13.200,00}} \\
    BOTUMEDIC IMPORTAÇÃO E EXPORTAÇÃO EIRELLI & \textcolor{accentColor}{\textbf{8.045,21}} \\
    GABRIELLA MACHADO COSTA LISBOA & \textcolor{accentColor}{\textbf{7.670,00}} \\
    TALMAX PRODUTOS DE PROTESE DENTARIA LTDA & \textcolor{accentColor}{\textbf{6.854,19}} \\
    55DENTAL CREMER PRODUTOS ODONTOLOGICOS S.A & \textcolor{accentColor}{\textbf{5.642,06}} \\
    NADIR PACIENTE & 2.940,00 \\
    SOUSAM IMPORTACAO E EXPORTACAO LTDA & 1.492,02 \\
    OUTROS (13 fornecedores) & 1.895,35 \\
    \midrule
    \rowcolor{lightGray} \textbf{TOTAL (24 credores)} & \textbf{629.317,28} \\
\end{longtable}

\vspace{-0.3cm}
\begin{quote}
    \footnotesize \textbf{NOTA VERSÃO 2:} Esta projeção ajustada exclui UNIMED GOIÂNIA e SIMONE HELENA DOS SANTOS, e considera SIMPLES NACIONAL reduzido a 2/3 devido à previsão de menor volume de notas fiscais. Total de 24 credores ativos, com forte concentração no Santander Brasil (80,9\%).
\end{quote}

\newpage

% --- SEÇÃO 3: TÍTULOS DE ALTA PRIORIDADE ---
\section{Títulos de Alta Prioridade}

Relação dos principais títulos com valor superior a \textbf{R\$ 5.000,00}, ordenados por valor decrescente. Estes títulos exigem atenção especial para planejamento de fluxo de caixa.

\begin{longtable}{c p{5.5cm} r p{3cm}}
    \toprule
    \rowcolor{tableHeader} \textbf{Vencimento} & \textbf{Credor} & \textbf{Valor (R\$)} & \textbf{Referência} \\
    \midrule
    \endfirsthead

    \multicolumn{4}{c}{{\tablename\ \thetable{} -- continuação}} \\
    \toprule
    \rowcolor{tableHeader} \textbf{Vencimento} & \textbf{Credor} & \textbf{Valor (R\$)} & \textbf{Referência} \\
    \midrule
    \endhead

    \midrule
    \multicolumn{4}{r}{\textit{Continua na próxima página...}} \\
    \endfoot

    \bottomrule
    \endlastfoot

    26/01/2026 & SANTANDER BRASIL & \textcolor{accentColor}{\textbf{11.292,42}} &  \\
    25/02/2026 & SANTANDER BRASIL & \textcolor{accentColor}{\textbf{7.363,11}} &  \\
    09/06/2031 & SANTANDER BRASIL & \textcolor{accentColor}{\textbf{7.300,01}} &  \\
    09/02/2031 & SANTANDER BRASIL & \textcolor{accentColor}{\textbf{7.300,01}} &  \\
    11/03/2031 & SANTANDER BRASIL & \textcolor{accentColor}{\textbf{7.300,01}} &  \\
    10/04/2031 & SANTANDER BRASIL & \textcolor{accentColor}{\textbf{7.300,01}} &  \\
    10/05/2031 & SANTANDER BRASIL & \textcolor{accentColor}{\textbf{7.300,01}} &  \\
    09/07/2031 & SANTANDER BRASIL & \textcolor{accentColor}{\textbf{7.300,01}} &  \\
    08/02/2030 & SANTANDER BRASIL & \textcolor{accentColor}{\textbf{7.300,00}} &  \\
    10/03/2030 & SANTANDER BRASIL & \textcolor{accentColor}{\textbf{7.300,00}} &  \\
    09/04/2030 & SANTANDER BRASIL & \textcolor{accentColor}{\textbf{7.300,00}} &  \\
    09/05/2030 & SANTANDER BRASIL & \textcolor{accentColor}{\textbf{7.300,00}} &  \\
    08/06/2030 & SANTANDER BRASIL & \textcolor{accentColor}{\textbf{7.300,00}} &  \\
    08/07/2030 & SANTANDER BRASIL & \textcolor{accentColor}{\textbf{7.300,00}} &  \\
    07/08/2030 & SANTANDER BRASIL & \textcolor{accentColor}{\textbf{7.300,00}} &  \\
    04/01/2031 & SANTANDER BRASIL & \textcolor{accentColor}{\textbf{7.300,00}} &  \\
    05/12/2030 & SANTANDER BRASIL & \textcolor{accentColor}{\textbf{7.300,00}} &  \\
    05/11/2030 & SANTANDER BRASIL & \textcolor{accentColor}{\textbf{7.300,00}} &  \\
    06/10/2030 & SANTANDER BRASIL & \textcolor{accentColor}{\textbf{7.300,00}} &  \\
    06/09/2030 & SANTANDER BRASIL & \textcolor{accentColor}{\textbf{7.300,00}} &  \\
\end{longtable}

\vspace{-0.3cm}
\begin{quote}
    \footnotesize \textbf{ATENÇÃO:} Esta tabela apresenta os 20 primeiros títulos de alta prioridade. O relatório completo identifica \textbf{68 títulos} acima de R\$ 5.000,00, representando 36,6\% do total de títulos e concentrados principalmente em financiamentos bancários de longo prazo (2027-2031).
\end{quote}

\newpage

% --- SEÇÃO 4: CRONOGRAMA DE PAGAMENTOS 2026 ---
\section{Cronograma de Pagamentos - Exercício 2026}

\subsection{Distribuição Mensal}

A tabela a seguir apresenta a projeção de desembolso mensal para o exercício de 2026, permitindo planejamento de liquidez.

\begin{table}[H]
    \centering
    \renewcommand{\arraystretch}{1.3}
    \begin{tabularx}{0.85\textwidth}{X r r}
        \toprule
        \rowcolor{tableHeader} \textbf{Mês/Ano} & \textbf{Valor (R\$)} & \textbf{\% do Total} \\
        \midrule
        Janeiro/2026 & 34.903,82 & 5,5\% \\
        Fevereiro/2026 & 34.700,40 & 5,5\% \\
        Março/2026 & 27.989,71 & 4,4\% \\
        Abril/2026 & 18.850,76 & 3,0\% \\
        Maio/2026 & 15.326,27 & 2,4\% \\
        Junho/2026 & 13.826,24 & 2,2\% \\
        Julho/2026 & 13.826,22 & 2,2\% \\
        Agosto/2026 & 13.926,44 & 2,2\% \\
        Setembro/2026 & 13.616,84 & 2,2\% \\
        Outubro/2026 & 13.616,84 & 2,2\% \\
        Novembro/2026 & 13.616,84 & 2,2\% \\
        Dezembro/2026 & 13.616,84 & 2,2\% \\
        \midrule
        \rowcolor{lightGray} \textbf{TOTAL 2026} & \textbf{227.817,22} & \textbf{36,2\%} \\
        \bottomrule
    \end{tabularx}
    \caption*{\footnotesize \textbf{Nota:} Os 63,8\% restantes (\reais{401.500,06}) referem-se a obrigações com vencimento em 2027 a 2031, concentradas em financiamentos bancários.}
\end{table}

\subsection{Visualização Gráfica}

O gráfico abaixo demonstra a concentração temporal das obrigações ao longo de 2026, evidenciando a necessidade crítica de liquidez no primeiro trimestre.

\begin{figure}[H]
    \centering
    \begin{tikzpicture}
        \begin{axis}[
            ybar,
            width=14cm,
            height=8cm,
            bar width=15pt,
            ylabel={Valor (R\$ milhares)},
            xlabel={Mês - 2026},
            xtick=data,
            xticklabels={Jan, Fev, Mar, Abr, Mai, Jun, Jul, Ago, Set, Out, Nov, Dez},
            x tick label style={font=\small},
            y tick label style={font=\small},
            ymin=0,
            ymax=50,
            ymajorgrids=true,
            grid style={dashed, gray!30},
            nodes near coords,
            nodes near coords align={vertical},
            every node near coord/.append style={font=\tiny, rotate=90, anchor=west},
            enlarge x limits=0.05,
        ]
        \addplot[fill=primaryColor, draw=primaryColor!80!black] coordinates {
            (1, 34.9)
            (2, 34.7)
            (3, 28.0)
            (4, 18.9)
            (5, 15.3)
            (6, 13.8)
            (7, 13.8)
            (8, 13.9)
            (9, 13.6)
            (10, 13.6)
            (11, 13.6)
            (12, 13.6)
        };
        \end{axis}
    \end{tikzpicture}
    \caption{Distribuição Mensal de Pagamentos - 2026}
\end{figure}

\vspace{0.5cm}

\textbf{Análise Crítica do Fluxo - VERSÃO 2:}
\begin{itemize}
    \item \textbf{Redução Significativa:} Versão 2 apresenta valores mensais entre R\$ 13,6 mil e R\$ 34,9 mil, substancialmente menores que V1.
    \item \textbf{Concentração no 1º Bimestre:} Janeiro e fevereiro mantêm valores similares (\textbf{R\$ 34,9k e R\$ 34,7k}), representando 11,0\% do total geral.
    \item \textbf{1º Trimestre:} Representa \textbf{R\$ 97,6 mil} (15,5\% do total geral), com redução de 17,0\% em relação à V1.
    \item \textbf{Média Mensal 2026:} Aproximadamente R\$ 19,0 mil/mês (31,4\% menor que V1), indicando alívio significativo no fluxo de caixa.
\end{itemize}

\newpage

% --- SEÇÃO 5: ANÁLISE DETALHADA DOS PRÓXIMOS 6 MESES ---
\section{Análise Detalhada dos Próximos 6 Meses (Jan-Jun/2026)}

Esta seção apresenta uma análise aprofundada das obrigações financeiras projetadas para o período de Janeiro a Junho de 2026, totalizando \textbf{R\$ 145.597,20} (23,1\% do total geral) na Versão 2.

\subsection{Distribuição Mensal Detalhada - 1º Semestre 2026}

\begin{table}[H]
    \centering
    \renewcommand{\arraystretch}{1.4}
    \begin{tabularx}{\textwidth}{X r r r}
        \toprule
        \rowcolor{tableHeader} \textbf{Mês/2026} & \textbf{Quantidade} & \textbf{Valor Total (R\$)} & \textbf{\% Acumulado} \\
        \midrule
        Janeiro & 28 títulos & 34.903,82 & 5,5\% \\
        Fevereiro & 23 títulos & 34.700,40 & 11,1\% \\
        Março & 20 títulos & 27.989,71 & 15,5\% \\
        Abril & 11 títulos & 18.850,76 & 18,5\% \\
        Maio & 9 títulos & 15.326,27 & 20,9\% \\
        Junho & 7 títulos & 13.826,24 & 23,1\% \\
        \midrule
        \rowcolor{lightGray} \textbf{TOTAL (6 meses)} & \textbf{98 títulos} & \textbf{145.597,20} & \textbf{23,1\%} \\
        \bottomrule
    \end{tabularx}
    \caption*{\footnotesize Valores V2 baseados em análise de vencimentos excluindo UNIMED, SIMONE HELENA e SIMPLES reduzido. Ticket médio de R\$ 1.486,00 no período.}
\end{table}

\subsection{Análise de Composição: Bancário vs Operacional}

\begin{table}[H]
    \centering
    \renewcommand{\arraystretch}{1.4}
    \begin{tabularx}{\textwidth}{X r r r r}
        \toprule
        \rowcolor{tableHeader} \textbf{Mês} & \textbf{Bancário (R\$)} & \textbf{Operacional (R\$)} & \textbf{Total (R\$)} & \textbf{\% Bancário} \\
        \midrule
        Janeiro & 11.292,42 & 23.611,40 & 34.903,82 & 32,3\% \\
        Fevereiro & 14.663,11 & 20.037,29 & 34.700,40 & 42,3\% \\
        Março & 9.866,52 & 18.123,19 & 27.989,71 & 35,3\% \\
        Abril & 9.026,88 & 9.823,88 & 18.850,76 & 47,9\% \\
        Maio & 7.904,58 & 7.421,69 & 15.326,27 & 51,6\% \\
        Junho & 7.904,58 & 5.921,66 & 13.826,24 & 57,2\% \\
        \midrule
        \rowcolor{lightGray} \textbf{TOTAL} & \textbf{60.658,09} & \textbf{84.939,11} & \textbf{145.597,20} & \textbf{41,7\%} \\
        \bottomrule
    \end{tabularx}
    \caption*{\footnotesize \textbf{Bancário:} Santander Brasil. \textbf{Operacional V2:} Fornecedores e obrigações tributárias (SIMPLES reduzido), excluindo UNIMED e SIMONE HELENA.}
\end{table}

\subsection{Indicadores de Liquidez Mensal}

\begin{table}[H]
    \centering
    \renewcommand{\arraystretch}{1.4}
    \begin{tabularx}{\textwidth}{X r r r}
        \toprule
        \rowcolor{tableHeader} \textbf{Mês/2026} & \textbf{Maior Título (R\$)} & \textbf{Ticket Médio (R\$)} & \textbf{Desembolso Diário Médio (R\$)} \\
        \midrule
        Janeiro & 11.292,42 & 1.246,56 & 1.126,25 \\
        Fevereiro & 7.363,11 & 1.388,02 & 1.239,30 \\
        Março & 7.300,01 & 1.272,71 & 902,86 \\
        Abril & 7.300,00 & 1.450,06 & 628,36 \\
        Maio & 7.300,00 & 1.393,30 & 494,46 \\
        Junho & 7.300,00 & 1.536,03 & 460,87 \\
        \midrule
        \rowcolor{lightGray} \textbf{MÉDIA SEMESTRAL} & \textbf{7.976,09} & \textbf{1.381,11} & \textbf{808,68} \\
        \bottomrule
    \end{tabularx}
    \caption*{\footnotesize Desembolso diário médio calculado com base em 31 dias para Jan/Mar/Mai e 30 dias para Fev/Abr/Jun.}
\end{table}

\newpage

\subsection{Visualização da Composição - 1º Semestre 2026}

\begin{figure}[H]
    \centering
    \begin{tikzpicture}
        \begin{axis}[
            ybar stacked,
            width=14cm,
            height=8cm,
            bar width=18pt,
            ylabel={Valor (R\$ milhares)},
            xlabel={Mês - 2026},
            xtick=data,
            xticklabels={Janeiro, Fevereiro, Março, Abril, Maio, Junho},
            x tick label style={font=\small, rotate=30, anchor=east},
            y tick label style={font=\small},
            ymin=0,
            ymax=140,
            ymajorgrids=true,
            grid style={dashed, gray!30},
            legend style={at={(0.5,-0.25)}, anchor=north, legend columns=2, font=\small},
            enlarge x limits=0.15,
        ]
        \addplot[fill=primaryColor, draw=primaryColor!80!black] coordinates {
            (1, 11.3) (2, 14.7) (3, 9.9) (4, 9.0) (5, 7.9) (6, 7.9)
        };
        \addplot[fill=accentColor!60, draw=accentColor!80!black] coordinates {
            (1, 23.6) (2, 20.0) (3, 18.1) (4, 9.8) (5, 7.5) (6, 6.0)
        };
        \legend{Obrigações Bancárias, Obrigações Operacionais}
        \end{axis}
    \end{tikzpicture}
    \caption{Composição de Pagamentos: Bancário vs Operacional (Janeiro a Junho 2026)}
\end{figure}

\subsection{Principais Credores no 1º Semestre 2026}

\begin{table}[H]
    \centering
    \renewcommand{\arraystretch}{1.3}
    \begin{tabularx}{\textwidth}{X r r r}
        \toprule
        \rowcolor{tableHeader} \textbf{Credor} & \textbf{Títulos} & \textbf{Valor (R\$)} & \textbf{\% do Semestre} \\
        \midrule
        SANTANDER BRASIL & 76 & 60.658,09 & 41,7\% \\
        SIMPLES NACIONAL & 12 & 20.802,08 & 14,3\% \\
        MR TRIPLAN CONTABILIDADE & 12 & 7.200,00 & 4,9\% \\
        GFIP SEFIP 8,40 GRF FGTS & 12 & 6.600,00 & 4,5\% \\
        VICTALAB F MANIPULACAO & 6 & 10.863,64 & 7,5\% \\
        BOTUMEDIC IMPORTAÇÃO & 4 & 4.545,21 & 3,1\% \\
        GABRIELLA MACHADO & 2 & 2.670,00 & 1,8\% \\
        TALMAX PRODUTOS DE PROTESE & 3 & 1.854,19 & 1,3\% \\
        EDUARDO BORGES NUNES & 6 & 16.200,00 & 11,1\% \\
        WANESSA DIAS VARANDA & 6 & 9.000,00 & 6,2\% \\
        OUTROS & 7 & 5.204,03 & 3,6\% \\
        \midrule
        \rowcolor{lightGray} \textbf{TOTAL} & \textbf{146} & \textbf{145.597,24} & \textbf{100,0\%} \\
        \bottomrule
    \end{tabularx}
    \caption*{\footnotesize Distribuição detalhada dos credores com vencimento no primeiro semestre de 2026.}
\end{table}

\subsection{Gráfico de Evolução do Ticket Médio}

\begin{figure}[H]
    \centering
    \begin{tikzpicture}
        \begin{axis}[
            width=14cm,
            height=7cm,
            xlabel={Mês - 2026},
            ylabel={Ticket Médio (R\$)},
            xtick=data,
            xticklabels={Jan, Fev, Mar, Abr, Mai, Jun},
            x tick label style={font=\small},
            y tick label style={font=\small},
            ymin=0,
            ymax=2800,
            ymajorgrids=true,
            grid style={dashed, gray!30},
            mark=*,
            mark size=3pt,
            line width=1.5pt,
            color=primaryColor,
        ]
        \addplot coordinates {
            (1, 1246.56) (2, 1388.02) (3, 1272.71) (4, 1450.06) (5, 1393.30) (6, 1536.03)
        };
        \end{axis}
    \end{tikzpicture}
    \caption{Evolução do Ticket Médio por Título - 1º Semestre 2026}
\end{figure}

\newpage

\subsection{Análise de Criticidade por Mês}

\begin{table}[H]
    \centering
    \renewcommand{\arraystretch}{1.4}
    \begin{tabularx}{\textwidth}{X r r c}
        \toprule
        \rowcolor{tableHeader} \textbf{Mês/2026} & \textbf{Títulos > R\$ 5k} & \textbf{Valor Crítico (R\$)} & \textbf{Nível de Risco} \\
        \midrule
        Janeiro & 1 & 11.292,42 & \cellcolor{orange!30} \textbf{MÉDIO} \\
        Fevereiro & 2 & 14.663,11 & \cellcolor{accentColor!30} \textbf{ALTO} \\
        Março & 3 & 16.566,53 & \cellcolor{accentColor!30} \textbf{ALTO} \\
        Abril & 2 & 14.426,88 & \cellcolor{accentColor!30} \textbf{ALTO} \\
        Maio & 2 & 14.304,58 & \cellcolor{accentColor!30} \textbf{ALTO} \\
        Junho & 2 & 14.304,58 & \cellcolor{accentColor!30} \textbf{ALTO} \\
        \midrule
        \rowcolor{lightGray} \textbf{TOTAL} & \textbf{12} & \textbf{85.558,10} & -- \\
        \bottomrule
    \end{tabularx}
    \caption*{\footnotesize \textbf{Critério de Risco:} ALTO (>25\% do valor mensal em títulos únicos >R\$ 5k), MÉDIO (10-25\%), BAIXO (<10\%).}
\end{table}

\subsection{Projeção de Fluxo de Caixa Acumulado}

\begin{figure}[H]
    \centering
    \begin{tikzpicture}
        \begin{axis}[
            width=14cm,
            height=8cm,
            xlabel={Mês - 2026},
            ylabel={Valor Acumulado (R\$ milhares)},
            xtick=data,
            xticklabels={Jan, Fev, Mar, Abr, Mai, Jun},
            x tick label style={font=\small},
            y tick label style={font=\small},
            ymin=0,
            ymax=190,
            ymajorgrids=true,
            grid style={dashed, gray!30},
            mark=*,
            mark size=3pt,
            line width=2pt,
            color=accentColor,
            area style,
            fill=primaryColor!20,
        ]
        \addplot+[fill] coordinates {
            (1, 34.9) (2, 69.6) (3, 97.6) (4, 116.5) (5, 131.8) (6, 145.6)
        };
        \end{axis}
    \end{tikzpicture}
    \caption{Fluxo de Caixa Acumulado - Obrigações do 1º Semestre 2026}
\end{figure}
\pagebreak

\subsection{Resumo Analítico do Período}

\textbf{Principais Conclusões - 1º Semestre 2026:}

\begin{itemize}
    \item \textbf{Diversificação de Obrigações:} Diferente do cenário anterior, o semestre apresenta distribuição equilibrada entre credores, com Santander representando apenas 41,7\% (vs 82,3\% anterior).

    \item \textbf{Compromissos Estruturais Reduzidos:} SIMPLES NACIONAL (14,3\%) representa as principais obrigações fiscais, com redução significativa de custos operacionais no cenário V2.

    \item \textbf{Inversão de Padrão:} Obrigações operacionais (58,3\%) superam bancárias (41,7\%), indicando perfil de despesas recorrentes e obrigações fiscais otimizadas.

    \item \textbf{Crescimento do Ticket Médio:} Evolução de R\$ 1.247 em janeiro para R\$ 1.536 em junho, refletindo consolidação de pagamentos mensais otimizados.

    \item \textbf{Gestão de Risco Otimizada:} 12 títulos críticos (>R\$ 5k) concentrados em Fev-Jun, com redução significativa de 40\% vs cenário V1.
\end{itemize}

\newpage

% --- SEÇÃO 6: CONCLUSÕES E RECOMENDAÇÕES ---
\section{Conclusões e Recomendações Técnicas}

\subsection{Principais Achados}

\begin{enumerate}
    \item \textbf{Exposição Financeira Otimizada:} O montante total de \textbf{R\$ 629,3 mil} (redução de 14,8\% vs V1) representa estrutura de obrigações otimizada, com \textbf{36,6\%} dos títulos classificados como de alta prioridade.

    \item \textbf{Diversificação de Riscos:} Redução da dependência bancária de 88,7\% para \textbf{80,9\%}, com foco em obrigações essenciais: SIMPLES NACIONAL (6,6\%) como principal obrigação operacional.

    \item \textbf{Perfil de Despesas Recorrentes:} Distribuição equilibrada em 2026 (\textbf{R\$ 227,8 mil} - 36,2\%) com média mensal de R\$ 19,0 mil, indicando estabilização operacional otimizada.

    \item \textbf{Obrigações de Longo Prazo:} Aproximadamente 63,8\% das obrigações (\reais{401,5 mil}) estendem-se até 2031, concentradas em financiamentos bancários estruturados.
\end{enumerate}

\subsection{Recomendações Estratégicas}

\begin{itemize}
    \item \textbf{Gestão de Obrigações Recorrentes:} Estabelecer provisões mensais para SIMPLES NACIONAL (R\$ 3,5k) e contabilidade (R\$ 1,2k), totalizando ~R\$ 4,7k/mês em obrigações fixas operacionais.

    \item \textbf{Monitoramento de Liquidez:} Manter reserva mínima de R\$ 40k para cobertura de oscilações mensais e obrigações imprevistas, considerando média mensal de R\$ 19,0k.

    \item \textbf{Revisão de Financiamentos Bancários:} Avaliar possibilidade de renegociação dos 156 títulos de longo prazo (2027-2031) do Santander, visando redução de taxas ou alongamento de prazos.

    \item \textbf{Otimização Tributária:} Analisar planejamento tributário para otimização do SIMPLES NACIONAL (R\$ 41,6k anuais no cenário V2), considerando enquadramento e alíquotas aplicáveis.

    \item \textbf{Dashboard Financeiro:} Implementar sistema de acompanhamento mensal automático para as 24 linhas de credores, com alertas para vencimentos >R\$ 5k.

    \item \textbf{Política de Fornecedores:} Negociar prazos diferenciados com VICTALAB (R\$ 16,3k) e BOTUMEDIC (R\$ 8k) para melhor distribuição de fluxo de caixa.
\end{itemize}

\vspace{2cm}

% \begin{center}
%     \rule{12cm}{0.5pt}

%     \vspace{1cm}

%     \textbf{DOCUMENTO TÉCNICO GERADO POR PROCESSAMENTO AUTOMATIZADO} \\
%     \vspace{0.3cm}
%     {\small Baseado em dados extraídos do sistema contábil em 07/01/2026} \\
%     {\small Processamento: Python 3 + Análise de Dados CSV} \\

%     \vspace{1cm}

%     \textit{Este relatório foi elaborado para fins de gestão financeira interna.} \\
%     \textit{Todos os valores são expressos em Reais (R\$) e referem-se à competência indicada.}
% \end{center}

\end{document}
